\chapter{Invariancia causal de las interfaces respecto de referencias externas}
\label{chap:independencia-prospectiva}
\index{independencia prospectiva}\index{grafo causal}
\index{protocolo ciego}\index{torres!completitud reconstructiva}

\section{Criterio de independencia causal respecto de referencias externas: grafo dirigido acíclico y fijación previa de la salida}
\label{sec:definiciones-independencia-prospectiva}

Sean \(\mathsf P\) las primitivas internas, \(\mathsf R\) las reglas,
lectores, calibres y normalizaciones declarados, \(G\) un generador,
\(o=G(\mathsf P,\mathsf R)\) su salida y \(y\) una referencia externa.

\begin{definition}[Interfaz prospectivamente independiente]
\label{def:interfaz-prospectivamente-independiente}
La interfaz es prospectivamente independiente de \(y\) si:
\begin{enumerate}
 \item \(y\notin\operatorname{Anc}(o)\);
 \item \(\mathsf P,\mathsf R,G\) y los controles negativos se congelan antes
 de abrir \(y\);
 \item la ejecución no tiene acceso directo ni indirecto a \(y\);
 \item \(o\) y su traza se sellan antes de calcular una comparación
 \(\operatorname{cmp}(o,y)\).
\end{enumerate}
\end{definition}

\begin{definition}[Independencia de selección]
\label{def:independencia-seleccion}
La ceguera de cada ejecución no basta si, tras producir candidatos
\(G_j(\mathsf P,\mathsf R_j)\), se publica
\[
 j_\ast(y)=\arg\min_j
 \operatorname{cmp}\!\left(G_j(\mathsf P,\mathsf R_j),y\right).
\]
La independencia de selección prohíbe que esta operación determine la salida
publicada.
\end{definition}

\begin{definition}[Predicción multivaluada]
\label{def:prediccion-multivaluada}
Si los descriptores físicos no metrológicos y las restricciones internas
admiten varias multisecciones, la salida prospectiva es el conjunto completo
\[
 \widehat{\mathfrak R}(d)
 =\{\mathcal M:\mathcal M\text{ satisface todas las restricciones
 preinscritas para }d\}.
\]
La no unicidad es una salida del procedimiento. Elegir después un elemento
por proximidad a una magnitud está prohibido.
\end{definition}

\section{Teorema de no interferencia}
\label{sec:teorema-no-interferencia}

\begin{definition}[Grafo causal dirigido acíclico fiel y completo]
\label{def:dag-fiel-completo}
Un grafo dirigido acíclico \(D\) es fiel si cada nodo representa la misma
función y las mismas entradas que la construcción citada. Es completo para
una salida si toda dependencia causalmente pertinente aparece como arista o
raíz declarada. Fidelidad y completitud son hipótesis semánticas; la
aciclicidad del conjunto de datos no las demuestra.
\end{definition}

\begin{theorem}[Independencia prospectiva condicionada al grafo causal]
\label{thm:no-interferencia-prospectiva}
Sea un sistema determinista de ecuaciones estructurales representado por un
grafo dirigido acíclico finito \(D\), fiel y completo para \(o\). Si
\(y\notin\operatorname{Anc}_D(o)\), todas las raíces ancestrales de \(o\)
permanecen congeladas y la ejecución no tiene acceso oculto a \(y\), entonces
para cualesquiera \(y_0,y_1\),
\[
 o_{\operatorname{do}(y=y_0)}
 =o_{\operatorname{do}(y=y_1)}.
\]
\end{theorem}

\begin{proof}
Se ordenan topológicamente los antecesores de \(o\). Sus raíces son internas
y permanecen fijas. Por inducción sobre el orden, cada nodo recibe los mismos
argumentos bajo ambas intervenciones y, por determinismo, devuelve el mismo
valor. La igualdad se propaga hasta \(o\). El resumen criptográfico no crea
la independencia; sólo hace detectable una alteración posterior.
\end{proof}

\begin{corollary}[Permutación de la comparación externa]
\label{cor:permutacion-comparacion-externa}
Si se permutan nombres, filas o valores externos después del sellado, la
salida interna no cambia, aunque pueda cambiar su puntuación:
\[
 o_{\operatorname{do}(y=\sigma y_0)}
 =o_{\operatorname{do}(y=y_0)},\qquad
 \operatorname{cmp}(o,\sigma y_0)
 \ne\operatorname{cmp}(o,y_0).
\]
\end{corollary}

\begin{proposition}[Separación de ramas calibradas]
\label{prop:separacion-ramas-calibradas}
Si una referencia externa intervino en una selección histórica, la
dependencia no contamina los objetos internos que no descienden de esa
selección. La rama puede aislarse y someterse a un protocolo prospectivo
nuevo sin modificar el núcleo.
\end{proposition}

\begin{proof}
Ser antecesor es una propiedad local del grafo. Las ramas
\[
 Z_0\longrightarrow p=5\longrightarrow\text{escala electrónica},
\qquad
 \{\text{masas}\}\longrightarrow\{\text{firmas calibradas}\}
 \longrightarrow\text{tabla}
\]
quedan fuera del subgrafo protegido. Las torres, el carácter formal, el
atlas y el centro \((5,5)\) no reciben aristas desde esas referencias.
\end{proof}

\section{Invariancia contrafactual sobre el grafo causal dirigido y alcance de la conclusión}
\label{sec:certificado-independencia-prospectiva}

El grafo causal publicado contiene \(71\) nodos, \(89\) aristas,
\(16\) interfaces y
\(29\) salidas protegidas. La comprobación causal establece aciclicidad, recuentos,
rutas declaradas, existencia de las fuentes registradas, ausencia de
referencias externas entre los antecesores consignados e invariancia
contrafactual del modelo simbólico. El resultado es
\[
 \#\{\text{antecesores externos de salidas protegidas}\}=0.
\]
El certificado no demuestra por sí solo que no falte una arista ni audita
todas las lecturas de datos de cada programa científico. Su fuerza es, por
tanto, \textsc{demostrado con estructura de partida explícita}: la conclusión
se obtiene sobre el grafo dirigido acíclico expresamente declarado y queda
condicionada a su fidelidad y completitud.

La sustitución simultánea de las referencias CODATA, Planck,
Barbero--Immirzi, CKM, masas, gravedad y cristal temporal puede alterar las
comparaciones y las ramas calibradas. No altera, dentro del grafo declarado,
el catálogo de \(468\) emisiones, las cinco regiones de \(\pi\), la
generación coinductiva correlacionada, \(\alpha_{12}\), \(C^\ast\), la década \(-34\), las torres,
el atlas, el centro electrónico, la doble proyección ni el corredor
excepcional.

\section{Ensayo multivaluado sin acceso al valor objetivo y criterio de identificación}
\label{sec:protocolo-ciego-multivaluado}

\begin{algorithm}[Asignación futura sin magnitudes]
\label{alg:asignacion-futura-sin-magnitudes}
La interfaz física se ejecuta en cuatro fases irreversibles:
\begin{enumerate}
 \item \emph{Congelación}: se sellan autoridad, atlas, involuciones,
 restricciones, código, prioridad y controles negativos.
 \item \emph{Entrada permitida}: sólo se reciben sector, generación, carga,
 color, conjugación, espín/estadística, composición y tipo metrológico; se
 excluyen masas, firmas refinadas, residuos, ángulos y unidades.
 \item \emph{Salida}: se enumeran todas las multisecciones compatibles; para
 cada extensión de cada multisección se calculan ruta observable, registro,
 firma, \(\chi_{\rm form}\), especialización y observable. En objetos
 compuestos se empalman primero las extensiones y los registros. Se sella el
 multiconjunto completo.
 \item \emph{Apertura}: un custodio distinto abre las magnitudes y calcula
 residuos, incertidumbres y covarianzas sin modificar candidato alguno.
\end{enumerate}
\end{algorithm}

La multisección enriquecida, y no una extensión puntual seleccionada
retrospectivamente, es la salida del emparejamiento físico. La evaluación
posterior puede ejecutarse fibra a fibra. El algoritmo precedente es una
especificación falsable; no se presenta como una predicción multiconjunto ya
ejecutada.

\begin{criterion}[Fallo del protocolo ciego]
\label{crit:fallo-protocolo-ciego}
El protocolo falla si el código puede leer una magnitud, si una firma
histórica calibrada reaparece como entrada, si se descarta un candidato tras
abrir la referencia, si se confunde masa positiva con masa exactamente nula
o si el multiconjunto sellado cambia al permutar el catálogo externo.
\end{criterion}

\section{Carácter formal y completitud reconstructiva}
\label{sec:completitud-reconstructiva-torres}

Para \(v=(n_A,s_C,k,b_{120},b_{\zeta})\in\ZZ^5\), el carácter formal
anterior a toda especialización es
\[
 \chi_{\rm form}(v;x_A,x_C,x_\Delta,x_{120},x_{270})
 =\exp\!\left(
 n_Ax_A+s_Cx_C+kx_\Delta
 +b_{120}x_{120}+b_{\zeta}x_{270}\right).
\]
La especialización declarada
\[
 \Theta_0
 =\left(A_{\rm rad},\frac{C^\ast_{\rm rad}}6,
 \Dfour,s_{120},\zeta_{270}\right),
 \qquad \zeta_{270}=135\Delta_4^2,
\]
produce
\(\chi_0(v)
=\chi_{\rm form}(v;\Theta_0)\).
La firma cruda, el carácter formal, su especialización central, el
refinamiento global \(\eta\in\mathcal E_{\mathfrak S}\) y una tabla
calibrada son cinco objetos diferentes.

\begin{theorem}[Completitud reconstructiva positiva]
\label{thm:completitud-reconstructiva-positiva}
Sea
\[
 \mathcal N=\langle90,120\rangle_{>0}
\]
y sea \((n_j)_{j\ge1}\) su enumeración creciente. Dados
\(0<q_+<q_-<1\), definamos
\[
 a_j=q_-^{n_j}-q_+^{n_j}>0.
\]
Entonces \(a_j\to0\) y \(\sum_ja_j<\infty\). Fijada una regla determinista
para resolver empates de semienteros, para todo \(r\in\RR\) la recurrencia
\[
 r_0=r,\qquad
 \nu_j=\operatorname{nint}\!\left(\frac{r_{j-1}}{a_j}\right),\qquad
 r_j=r_{j-1}-\nu_ja_j
\]
satisface
\[
 |r_j|\leq\frac{a_j}{2}\longrightarrow0,\qquad
 r=\sum_{j\ge1}\nu_ja_j,\qquad
 \sum_{j\ge1}|\nu_ja_j|<\infty.
\]
Por tanto, todo \(R>0\), con \(r=\log R\), se representa exactamente en la
completación por
\(\exp(\sum_j\nu_ja_j)\).
\end{theorem}

\begin{proof}
Los índices de \(\mathcal N\) son \(90,120\) y todos los múltiplos de \(30\)
desde \(180\). En consecuencia,
\[
 \sum_{n\in\mathcal N}q_-^n
 =q_-^{90}+q_-^{120}
 +\frac{q_-^{180}}{1-q_-^{30}}<\infty.
\]
Como \(0<a_j<q_-^{n_j}\), la sucesión es sumable y converge a cero. La
definición de entero más próximo implica
\(|r_j|\le a_j/2\). Por telescopado,
\(\sum_{i=1}^{N}\nu_ia_i=r-r_N\to r\). Para \(j\ge2\),
\[
 |\nu_ja_j|=|r_{j-1}-r_j|
 \le\frac{a_{j-1}+a_j}{2},
\]
y la sumabilidad de \((a_j)\) prueba la convergencia absoluta.
\end{proof}

El teorema es una completitud del diccionario sobre \(\RR_{>0}\). Si los
coeficientes \(\nu_j\) se obtienen del residuo de una masa conocida, su uso
es \textsc{reconstrucción calibrada}; no proporciona la selección
prospectiva. Tampoco representa masa exactamente nula. Fotón, gluón o un
eventual modo de espín dos requieren protección exacta por simetría en la
representación física.

\Needspace{34\baselineskip}
\section{Protocolos prospectivos de invariancia para las familias de constantes}
\label{sec:protocolos-prospectivos-especificos}

\begin{description}
 \item[Barbero--Immirzi.] Se congelan \(S_{90},S_{120}\), los pesos
 \((12,-1)\) y su normalización antes de abrir una referencia física. La
 evaluación de \(\Gamma_{6\to8}\) es ciega respecto de esa referencia,
 pero la unicidad previa de series y pesos y el entrelazador con el operador
 de área son problemas distintos.
 \item[CKM y CP.] Se sellan las cuatro ecuaciones, sus coeficientes y la
 convención angular; se publican simultáneamente los cuatro parámetros,
 \(|V|\) y \(J\). El contraste usa una covarianza fijada de antemano.
 \item[Constitutivas.] Las identidades adimensionales preceden a
 \(Z_{\rm ref}\) y \(c_{\rm ref}\). El cambio de carta transforma la
 realización, no las identidades, y no crea cuatro predicciones
 independientes.
 \item[Gravedad e inercia.] Las razones
 \(G_{\rm tor}/I_{\rm tor}=1\) y
 \(G_{\rm av}/I_{\rm av}=\pi/\varphi^2\) pertenecen a capas diferentes.
 Su interpretación física exige un morfismo y un contraste prospectivo; un
 fenómeno cosmológico no selecciona la razón.
\end{description}

\section{Estratificación probatoria: identidades internas, deducciones, reconstrucciones calibradas e interfaces sometidas a cierre matemático}
\label{sec:cinco-niveles-independencia}

\begin{table}[htbp]
\centering
\caption{Niveles que la independencia causal no permite colapsar.}
\label{tab:cinco-niveles-independencia}
\small
\begin{tabularx}{\textwidth}{@{}L{0.25\textwidth}YL{0.27\textwidth}@{}}
\toprule
Nivel & Ejemplo & Afirmación permitida\\
\midrule
Generador interno certificado&
Regiones, \(\alpha_{12}\), \(C^\ast\), \(-34\), atlas&
Salida obtenida en el dominio declarado.\\
Evaluación condicionada ciega&
\(\Gamma_{6\to8}\), CKM fijado, \(\chi_0(v)\)&
No consulta el blanco, pero depende del funcional o firma.\\
Protocolo ciego no ejecutado&
Especificación multivaluada&
Existe un procedimiento falsable; no una salida sellada.\\
Reconstrucción calibrada&
Firmas históricas y refinamiento desde \(r=\log R\)&
Evaluación exacta condicionada a una referencia.\\
Interfaz física pendiente&
Partícula--antipartícula, área, masa nula, gravedad&
Dominio definido; falta representación o contraste.\\
\bottomrule
\end{tabularx}
\end{table}
